\documentclass[8pt]{strategynet-journal}
\usepackage{volt-paper-style}
\hypersetup{pdftitle={Learning to Hedge with Stocks and Options},pdfauthor={research@strategynet.ai}}
\newcommand{\E}{\mathbb{E}}
\SNHeader{strategynet.ai:Volt and Jormungandr}
\title{Learning to Hedge with Stocks and Options\\\large The Control Problem, Learning Method and Application Interfaces}
\author{\href{mailto:research@strategynet.ai}{research@strategynet.ai}}
\date{October 2026 --- Methods and implementation report}
\begin{document}
\maketitle
\begin{abstract}
\normalsize A dealer who sells an option must finance its eventual payment while managing the risk of market moves along the way. Trading stock offsets some of that exposure; adding a traded option gives the dealer another way to manage curvature and volatility risk. Every adjustment also incurs costs. We formulate this decision as a finite-horizon stochastic control problem, specifying the market dynamics, observations, holdings, cash account and terminal loss. We explain how squared loss and conditional value at risk lead to different learning objectives, and how a neural policy learns from complete trading paths through automatic differentiation. Volt supplies the financial calculation and learner, while Jormungandr supplies the training and inference interfaces. A worked cash account, a dated decision timeline and a tail-loss diagram connect the mathematics to an application. Numerical comparisons and a synthetic learning exercise demonstrate the calculation and its application interface. The client workflow covers scenario submission, training, held-out evaluation, holdings inference and checkpoint recovery.
\end{abstract}
\begin{multicols}{2}
\tableofcontents
\end{multicols}
\thispagestyle{fancy}
\clearpage
\begin{SNcolumns}
\section{Hedging an option book}
A dealer sells a call on a stock and receives its premium. At expiry, the
buyer can claim the excess of the stock price over the agreed strike. A put
pays the excess of the strike over the stock price. Selling both at the same
strike creates a short straddle: the dealer owes a payment when the stock
finishes far enough in either direction. The premium finances the hedge,
but the eventual payment and the trading needed to manage it remain uncertain.

A stock position offsets the option book's response to a small price change.
Delta measures that response, in shares per option unit. Gamma measures how
delta changes as the stock moves. Vega measures the response to a change in
volatility, with its quoted unit specified by the application. These exposures
change during the life of the book. A trader who continually adjusts a stock
hedge therefore pays transaction costs while retaining exposure to movements
between trading times and to changes in option prices.

A second option adds another way to manage that risk. Let $L$ denote the
liability and $C$ a traded call used for hedging. Given their model Greeks,
a delta--gamma hedge holds
\begin{equation}
 h^C=\frac{\Gamma_L}{\Gamma_C},\qquad
 h^S=\Delta_L-h^C\Delta_C.
\end{equation}
The call offsets curvature and the stock completes the delta adjustment.
The remaining vega is $\nu_L-h^C\nu_C$; its size depends on the chosen
contracts. A small hedge-option gamma can also produce a very large call
position. Instrument choice, position limits and trading costs matter as much
as the local sensitivity calculation.

Deep hedging learns the holdings from simulated trading experience~\cite{deep}.
Fran\c{c}ois and colleagues develop the stock-and-option setting using joint
index returns and implied-volatility information~\cite{options,reference}.
The policy receives current observations and existing positions, then selects
the next stock and option quantities. Training rewards decisions that reduce
a specified measure of the final hedging loss after costs and financing.
This is useful when a simple Greek-matching rule trades too frequently, or
when several instruments offer different combinations of risk reduction and
execution cost. Model hedges still provide valuable comparisons: they explain
what the neural policy gains or gives up as the market and cost assumptions change.

\section{The market and the decision process}
Write the trading dates as $t_k=k\Delta t$, $k=0,\ldots,N$. The episode
ends at $T=t_N$. The policy chooses holdings at $k=0,\ldots,N-1$; the final
date settles the liability and sells the hedge. Daily trading means one action
per trading day. Simulation may use finer time steps to construct prices
between those dates while preserving the chosen trading frequency.
\end{SNcolumns}

\begin{center}
\begin{minipage}{\textwidth}
\noindent Figure~\ref{fig:episode} follows one trading episode. Each decision uses the
observations and holdings available at that date; the final loss supplies the
training signal for all preceding actions.
\par\medskip\centering
\begin{tikzpicture}[>=Latex,font=\normalsize]
\draw[->,thick,darkblue] (0,0)--(13,0);
\foreach \x/\lab in {0.5/{0},4.4/{t},8.3/{t+1},12.1/{T}}{
 \fill[darkblue] (\x,0) circle(2pt);
 \node[above=5pt] at(\x,0) {$\lab$};}
\node[align=center,text width=2.6cm,below=8pt] at(0.5,0) {Collect premium\\Choose holdings};
\node[align=center,text width=2.9cm,below=8pt] at(4.4,0) {Observe prices\\and positions\\Trade; pay costs};
\node[align=center,text width=2.8cm,below=8pt] at(8.3,0) {Observe next market\\Update financed wealth};
\node[align=center,text width=2.5cm,below=8pt] at(12.1,0) {Sell hedges\\Pay liability\\Measure loss};
\draw[<->,darkgreen] (4.4,0.85)--node[above]{One decision interval $\Delta t$}(8.3,0.85);
\draw[->,darkgreen] (12.1,-1.6)--(12.1,-2.1)--node[below]{Training: terminal risk supplies derivatives for all earlier decisions}(0.5,-2.1)--(0.5,-1.6);
\end{tikzpicture}
\captionof{figure}{One episode follows the liability from sale to expiry. The policy acts at specified dates using information available at that date. Prices move between decisions; financing and trading costs enter the account before the terminal objective is calculated. Training repeats this calculation across many market paths.}\label{fig:episode}
\end{minipage}
\end{center}

\begin{SNcolumns}

A market generator describes how the external market state $X_k$ changes:
\begin{equation}
 X_{k+1}=F_k(X_k,\varepsilon_{k+1}),\qquad
 \varepsilon_{k+1}\sim P_k(\,\cdot\mid X_k).
\end{equation}
The scenario probability law determines which outcomes the objective weights.
A risk-neutral law supports a pricing-model benchmark; an estimated real-world
law supports a hedge assessed under that return distribution. Both require
consistent instrument marks and financing assumptions.
Here $X$ can contain the stock price, volatility variables, interest rates
and any additional factors required by the model. The current examples treat
the trader as a price taker: changing the hedge affects wealth and costs,
while the external price process follows its own law. Market impact would
require a transition model with the action among its inputs.

The two-day application control uses geometric Brownian motion,
\begin{equation}
 S_{k+1}=S_k\exp\!\left[(r-q-\tfrac12\sigma^2)\Delta t
                    +\sigma\sqrt{\Delta t}\,\varepsilon_{k+1}\right],
\end{equation}
with independent standard-normal innovations, $\Delta t=1/252$,
$S_0=100$, $\sigma=0.20$, $r=0.03$ and $q=0.01$.
The liability pays $|S_2-100|$. The hedge call has strike 100 and an initial
84-trading-day maturity; its mark at each date is its Black--Scholes value
using the remaining maturity and those same parameters.
These choices completely specify this small market experiment.

The separate Heston study lets instantaneous variance $v$ change randomly:
\begin{align}
 dS_t&=\mu S_t\,dt+S_t\sqrt{v_t}\,dB_t,\\
 dv_t&=\kappa(\bar v-v_t)\,dt+\xi\sqrt{v_t}\,dW_t,\\
 d\langle B,W\rangle_t&=\rho\,dt.
\end{align}
Mean reversion has speed $\kappa$ and long-run variance $\bar v$; $\xi$
controls fluctuations in variance and $\rho$ connects the two shocks.
Our zero-rate Heston comparison uses $\mu=0$, $S_0=100$,
$v_0=\bar v=0.04$, $\kappa=1$, $\xi=2$, $\rho=-0.7$ and 30 daily
intervals of $1/365$. Its second hedge instrument is a funded variance claim.
That experiment and the two-day stock-and-call control have different market
and instrument definitions. Each has its own scenario and result records.

The full economic state is $(k,X_k,V_k,h_{k-1},u_k)$: date, market variables,
account wealth, existing holdings and any running contract statistic $u_k$.
An Asian payoff, for example, requires a running average. The agent receives
an observation vector $x_k=\phi_k(X_{0:k},u_k)$ and its previous holdings.
The application chooses which quantities to expose. If a latent volatility
factor is omitted, the observed problem can require memory to infer it.
Wealth can also be supplied as a feature when the policy should react to it.
The supplied two-day policy observes scaled log spot and time remaining,
plus previous holdings.

The action $h_k=(h_k^S,h_k^C)\in\mathbb R^2$ specifies the desired stock
and option holdings. Orders are the changes $h_k-h_{k-1}$. The current
continuous-output native policy permits signed quantities; adding position
limits requires a specified admissible action set and its implementation.
The initial holdings are zero and initial wealth is the premium. The account
recurrence below, together with $F_k$ and the terminal payoff, defines the
controlled system completely.

\section{Stock, option and cash positions}
The account starts with wealth $V_0$, typically the premium collected for the
liability. At each trading time $t$, the policy chooses stock shares $h_t^S$
and hedge-option units $h_t^C$. Let $H_t=(S_t,C_t)$ be their prices, and let
$h_{-1}=0$. With proportional trading fees $\kappa_i$, the cash left after
rebalancing is
\begin{equation}
 c_t=V_t-h_t^S S_t-h_t^C C_t
       -\sum_{i\in\{S,C\}}\kappa_i H_t^i|h_t^i-h_{t-1}^i|.
\end{equation}
This cash balance earns or pays interest. Under a stock dividend-reinvestment
convention, the next pre-trade wealth is
\begin{equation}
 V_{t+1}=c_t R_t+h_t^S S_{t+1}G_t+h_t^C C_{t+1},
\end{equation}
where $R_t=\exp(r_t\Delta t)$ and $G_t=\exp(q_t\Delta t)$. The scenario input
supplies these growth factors, allowing the account convention to be checked
alongside the prices. A negative cash balance represents borrowing.

At the liability's expiry, the remaining hedge positions are sold. With
closing-cost rates $\kappa_i^{\rm close}$, the available wealth and loss are
\begin{align}
 V_T^{\rm liq}&=V_T-\sum_i\kappa_i^{\rm close}H_T^i|h_{T-1}^i|,\\
 E_T&=L_T-V_T^{\rm liq}.
\end{align}
A positive $E_T$ is a shortfall; the seller's profit is $-E_T$.
The longer-dated hedge option is sold at its expiry-date market value while
the liability pays its contractual amount. Each opening, adjustment and
closing fee enters the account once.

Consider a one-period example. A dealer sells one call with strike 95 for a
premium of 6. The stock costs 100 and a longer-dated hedge call costs 4.
Buying 0.4 shares and 0.5 hedge-option units costs 42. Stock fees of 5 basis
points and option fees of 1\% add 0.04, leaving cash of $-36.04$. Over the
period, that borrowing grows by a factor of 1.001. The stock then reaches
102 and the hedge call 4.8; dividends are zero.

\begin{center}
\begin{tikzpicture}[font=\normalsize,x=1cm,y=1cm]
\draw[->] (0,0)--(5.7,0) node[right]{Gain};
\draw[->] (0,0)--(-1.15,0) node[left]{Cost};
\foreach \y/\len/\lab/\val in {0.7/3.6/{Stock gain}/{+0.80000},1.4/1.8/{Call gain}/{+0.40000}}{
 \fill[darkgreen!55] (0,\y-0.12) rectangle(\len,\y+0.12);
 \node[anchor=west] at(\len+0.12,\y) {\val};
 \node[anchor=west] at(0,\y+0.29) {\lab};}
\fill[darkblue!55] (-0.38,2.1-0.12) rectangle(0,2.1+0.12);
\node[anchor=west] at(0.15,2.1) {Fees: $-0.08440$};
\fill[darkblue!55] (-0.162,2.8-0.12) rectangle(0,2.8+0.12);
\node[anchor=west] at(0.15,2.8) {Financing: $-0.03604$};
\node[align=center] at(2.3,-0.65) {Premium minus liability: $6-7=-1$\\Total profit: $0.8+0.4-0.0844-0.03604-1=0.07956$};
\end{tikzpicture}
\captionof{figure}{Sources of profit in the two-date example. Bar lengths show currency contributions from trading gains, fees and financing. The premium and liability add a net cost of one currency unit.}
\end{center}

The profit can also be reconstructed from stock gains of 0.8, option gains
of 0.4, the original premium of 6, the liability payment of 7, total fees
of 0.0844 and financing of 0.03604. Both calculations give 0.07956.
This independent reconstruction helps catch sign and timing errors.

Contract multipliers must be applied consistently. An exchange contract
covering 100 shares corresponds to 100 option units when option prices are
quoted per share. The stock quantity is already in shares. The application
converts contracts to these economic quantities once, before constructing
the account. The same convention carries through premiums, Greeks, fees and
terminal payments. Scaling the entire example by 100 scales every cash flow
and the profit by 100; mean squared loss then scales by $100^2$.

\section{Learning from completed trading paths}
The policy is a function $h_t=\pi_\theta(x_t,h_{t-1})$, where $x_t$ contains
information available at the trading time and $\theta$ contains learned
weights. Spot, time to expiry, option quotes and volatility descriptors can
all be inputs. Existing holdings tell the policy how much its next decision
will cost to implement. Features drawn from a full future path must be
restricted to observations available at $t$.

For $B$ scenario paths, one objective is mean squared hedging loss,
\begin{equation}
 J(\theta)=\frac1B\sum_{b=1}^B E_T^{(b)}(\theta)^2.
\end{equation}
This penalizes both overfunding and shortfall. Downside squared loss replaces
$E_T^2$ by $\max(E_T,0)^2$, emphasizing the seller's unpaid obligation.
Conditional value at risk, or CVaR, concentrates on the worst losses. Its
level-$\alpha$ representation is
\begin{equation}
 \operatorname{CVaR}_\alpha(E_T)=\min_\eta
 \left\{\eta+\frac{\E[(E_T-\eta)_+]}{1-\alpha}\right\}.
\end{equation}
In reinforcement-learning notation, intermediate rewards can be set to zero
and the final reward to $-\ell(E_T)$, with discount factor one. For squared
loss, $\ell(e)=e^2$. For CVaR at a given threshold $\eta$,
\begin{equation}
 \ell_\eta(e)=\eta+\frac{(e-\eta)_+}{1-\alpha}.
\end{equation}
Optimizing $\eta$ alongside the policy yields the CVaR objective; the native
empirical calculation evaluates the equivalent fractional-tail statistic.
Funding already appears in the cash account. The trader commits at the start
to the chosen terminal-risk objective and learns a policy for that objective
across the full horizon. This formulation gives each action credit through
its eventual effect on the account, including future trading costs.

\begin{center}
\begin{tikzpicture}[>=Latex,font=\normalsize,x=1cm,y=1cm]
\draw[->] (0,0)--(6.3,0) node[right]{$e$};
\draw[->] (0,0)--(0,2.5) node[above]{$\ell_\eta(e)$};
\draw[thick,darkblue] (0.2,0.45)--(2.6,0.45)--(5.2,2.2);
\draw[dashed] (2.6,0)--(2.6,0.45);
\node[below] at(2.6,0) {$\eta$};
\node[anchor=west] at(0.3,0.75) {Constant below $\eta$};
\node[align=left] at(4.7,1.0) {Slope $1/(1-\alpha)$\\above $\eta$};
\end{tikzpicture}
\captionof{figure}{CVaR training puts additional weight on losses above the threshold. Increasing $\alpha$ concentrates the objective on a smaller, more adverse part of the loss distribution.}
\end{center}

The native learner supports squared, downside squared, empirical CVaR and
entropic objectives. Its empirical CVaR calculation weights a fractional
observation at the tail cutoff when the requested tail size is nonintegral.
Choice of objective changes the trading problem: a policy trained on average
squared error can differ materially from one protecting against extreme shortfall.

The terminal liability and trading account determine the training signal.
The learner runs each path forward, calculates the final loss, and propagates
derivatives back through all holdings, fees and cash balances to the weights.
An optimizer then changes those weights. Every decision date receives its learning signal from the final loss.
Known model hedges provide independent performance comparisons. The supplied continuous-action application uses
this direct differentiation approach within Jormungandr's learner interface.

Burn carries out automatic differentiation through the tensor calculation~\cite{burn}.
For exogenous simulated prices, those prices are fixed inputs during a weight
update; derivatives follow the policy and the account. Absolute-value trading
costs are piecewise differentiable, with the implementation's derivative
convention at zero turnover. Price Greeks use differentiation with respect
to financial inputs; policy training uses differentiation with respect to
network weights. Volt provides both applications, with distinct inputs and outputs.

For a batch of trajectories $\omega_b$, the actual update is
\begin{equation}
 \widehat J(\theta)=\frac1B\sum_b\ell(E_T(\theta;\omega_b)),\qquad
 \theta\leftarrow\operatorname{Adam}(\theta,\nabla_\theta\widehat J).
\end{equation}
The implemented method is direct policy optimization through a differentiable
simulator/account calculation. It uses the known derivatives of that calculation
to learn continuous actions. Other RL methods can estimate a value function
or a policy gradient from sampled rewards; integrating one of those engines
would retain the same economic problem specification. A policy's output is
the next holding, and its training target is the risk of the resulting account.

The current native policy uses a separate two-hidden-layer ReLU network at
each decision date and feeds previous holdings into the next decision.
It learns for the liability and scenario distribution supplied during training.
At deployment, changing market observations changes the action through
$x_t$. Supporting a broader range of books requires liability descriptors
and representative portfolios in the training distribution, followed by
validation across those books. Equivalent contractual decompositions should
produce the same net cash flows and economic exposures under matched costs;
that property can be tested directly in the account and policy inputs.

\end{SNcolumns}
\clearpage
\begin{SNcolumns}
\section{Training and inference through the API}
A hedging application supplies its scenarios and liability, then uses the
same model for training, evaluation and holdings inference. Volt evaluates
the financed account and differentiates its chosen risk objective.
Jormungandr manages the model and exposes the requests over HTTP~\cite{volt,jorm}.
A registered worker retains the policy, optimizer and scenario tensors
between requests. CPU, Metal and CUDA builds select where that calculation runs.

The following Python calls use \texttt{JormungandrClient}. Start the service
with \texttt{VOLT\_LEARNER\_BIN} pointing to the compiled Volt worker.
The public example \texttt{examples/native\_hedging.py} launches two
scenario-producing actors against it~\cite{jorm}. Here we follow one model
through the same API, using the two-day market defined in Section~2.



Here \texttt{premium} is the initial straddle value under the stated pricing
model. The default policy has two decision dates, two market features, two
hedge instruments and hidden width eight. Setting the objective to
\texttt{\{"kind":"cvar","alpha":0.95\}} selects 95\% CVaR; changing the
risk preference leaves the scenario submission workflow intact.

A batch supplies flattened arrays: \texttt{marks} has shape
$[B,N+1,2]$, \texttt{features} has $[B,N,2]$ and \texttt{payoffs} has $[B]$.
For each path the marks alternate stock and call prices at successive dates.
Features contain scaled log spot and time remaining at each decision.
The array \texttt{payoffs} holds the straddle payment owed at expiry.
Each batch also carries its path count and dataset identity with its first
path index. Registering disjoint ranges assigns the observations to training
and validation before the tensors are uploaded.





Training consumes the uploaded training batch and advances the model version
once per accepted update. Evaluation measures the current policy on the
separate validation batch. Replacing or streaming training batches extends
this loop to larger scenario sets. Multiple actors can submit work to one
model using the same versioned update protocol. The supplied example
implements retries after an actor receives a stale-version response.

For a current market observation, inference returns target stock and option
holdings. With spot 100, normalized log spot zero and the full two-day
horizon remaining, the input uses features $[0,1]$ and zero previous holdings.



The returned \texttt{holdings} vector uses the stock and call ordering of the
scenario marks. The application subtracts its current holdings to form orders.
A checkpoint preserves model weights, Adam moments, configuration and dataset
declarations. Restoring it reproduces holdings inference; continued training
re-uploads the required scenario tensors. The client releases resources with
\texttt{client.delete\_native\_model(model)}.

These calls map directly to HTTP: model creation uses
\texttt{POST /v2/native/models}; the other operations use
\texttt{POST /v2/native/models/\{id\}/\{operation\}}. Applications can use
these routes from another language or register another compatible policy engine.
The worker keeps the financial tensors and optimizer in its execution backend;
the application receives holdings, risk summaries and versioned model state.

\end{SNcolumns}
\par\noindent
\begin{minipage}{\textwidth}
\noindent The request sequence below creates the learner, submits two scenario batches,
trains it, evaluates held-out paths, retrieves holdings and restores a checkpoint.
The premium and scenario batches are the inputs described above.
\par\medskip
\begin{lstlisting}[basicstyle=\ttfamily\normalsize,columns=fullflexible,breaklines=true,frame=single,framesep=3pt]
import math
from jormungandr.client import JormungandrClient
client = JormungandrClient("http://127.0.0.1:8080")
model = "hedge"
client.create_native_model({
  "model_id": model, "backend": "volt_burn",
  "config": {
    "target": "cpu", "precision": "f64",
    "premium": premium,
    "objective": {"kind": "squared_loss"},
    "trading_costs": [0.0005, 0.01],
    "liquidation_costs": [0.0005, 0.01],
    "cash_growth": [math.exp(0.03/252)] * 2,
    "income_growth": [math.exp(0.01/252), 1.] * 2}})
client.native_command(model, "register_dataset", {
  "dataset_id": "gbm-study", "partitions": [
    {"split": "train", "first_path": 0, "paths": 64},
    {"split": "validation", "first_path": 64,
     "paths": 64}]})
for name, batch in [("train", training_batch),
                    ("validation", validation_batch)]:
    client.native_command(model, "upload", {
      "batch_id": name, "split": name, "batch": batch})
for _ in range(30):
    version = client.get_native_model(model)["version"]
    client.native_command(model, "update", {
      "batch_id": "train", "expected_version": version})
report = client.native_command(model, "evaluate", {
  "batch_id": "validation"})
targets = client.native_command(model, "infer", {
  "step": 0, "paths": 1, "features": [0., 1.],
  "previous_holdings": [0., 0.]})
state = client.native_command(model, "checkpoint")
state.pop("ok", None)
client.create_native_model({"model_id": "resumed",
  "backend": "volt_burn", "checkpoint": state})
\end{lstlisting}
\end{minipage}
\clearpage
\begin{SNcolumns}
\section{Validation and measured results}
We checked the trading calculation against a second implementation that
follows each market path one date at a time using ordinary numbers and loops.
It pays for holdings and fees, accrues interest and settles the liability.
Both versions were written for this project; their separate code paths help
catch sign errors, missing charges and incorrect payment timing.
Parameter derivatives were also checked by perturbing weights and comparing
the resulting change in loss.

The CPU double-precision checks agree to roughly $10^{-15}$ across the loss,
parameter derivatives and first optimizer update. The corresponding
CPU/Metal single-precision checks agree to roughly $10^{-6}$ for path profits.
The validation records retain the individual errors and test inputs.

A companion comparison also executed the original Deep Hedging author's
published financial engine and network components~\cite{authorcode}.
With matched inputs, both implementations reached validation CVaR of
approximately 0.262 in a two-date Black--Scholes learning control after
1,000 updates.

The end-to-end synthetic exercise used two trading days, a constant-volatility
stock process with annual volatility 20\%, interest 3\% and dividend yield
1\%. The liability was an at-the-money straddle and the hedge instruments
were stock and an initially 84-day at-the-money call. Black--Scholes prices
supplied the initial premium and hedge-call marks. Both networks had hidden
width eight. Trading fees were 5 basis points for stock and 1\% for the option,
including terminal liquidation.

Two actors each supplied one batch of 64 training paths and applied 30 updates,
for 60 accepted updates in total. Each actor reused its batch. On a separate
64-path validation set, mean squared hedging loss fell from 1.54 to
0.88, a 43\% reduction. This run exercises data submission, shared-model
updates and out-of-batch evaluation. Estimating economic performance calls
for larger independent samples, repeated training runs and direct comparisons
with model hedges.

\begin{center}
\begin{minipage}{\linewidth}
\centering
\begin{tabular}{@{}p{4.0cm}rr@{}}
\toprule
CPU workload & Volt & PyTorch\\
\midrule
64 paths, 2 dates, width 8 & 0.355 & 0.360\\
256 paths, 30 dates, width 17 & 5.957 & 5.440\\
\bottomrule
\end{tabular}
\captionof{table}{Median warm training-update time in milliseconds on an Apple M3 Max, macOS, float64, one CPU thread per runtime. Each update includes the trading calculation, backward differentiation and Adam optimization.}
\end{minipage}
\end{center}

The timing comparison used five warm-up updates and 20 alternating measured
updates. Volt used a tuned release build with ThinLTO and native CPU targeting.
Its time includes worker communication; PyTorch runs in process. Initial
upload and model creation are excluded. These measured workloads are small;
the table provides an execution reference for the application calculation.

Application checks also cover dataset-partition enforcement, stale model
versions, checkpoint continuation and a different registered learner engine.
The completed release validation included 288 Volt checks, 216 Jormungandr
checks and six native protocol/HTTP checks. These support the availability
of the application interface and numerical engine. The controls support using the API to construct and compare hedging policies.

\section{Using the workflow in a trading application}
The same interfaces support a research comparison and a deployed decision
process. For research, the application uploads common held-out market paths
and evaluates the learned policy alongside delta and delta--gamma hedges,
using the same premiums, instrument units, financing and execution costs.
Loss, turnover and account paths then explain the economic effect of the
policy. The mathematical market specification stays with the dataset so
that the comparison can be reconstructed.

For ongoing decisions, the application prepares the observation using the
training-time feature definitions, supplies current holdings and requests the
next target. The step index selects the policy for the remaining horizon.
Stock and option prices, contract maturity and any volatility features must
use the same conventions as training. Portfolio quantities convert the targets
into economic exposures; the execution application handles order sizes,
liquidity and position controls.

Changing the liability or market model begins with a new scenario specification.
Longer horizons, additional instruments and new market observations enter through
the configured policy dimensions and batch arrays. Predictive factors can enter
as features when their observation times and training distribution are defined.
This gives practitioners a common way to study whether additional information
or a different hedge instrument improves risk after trading costs.
\end{SNcolumns}
\clearpage
\begin{SNcolumns}
\begin{thebibliography}{99}
\interlinepenalty=10000
\bibitem{deep} H. Buehler, L. Gonon, J. Teichmann and B. Wood. \emph{Deep Hedging}. 2018. \url{https://arxiv.org/abs/1802.03042}.
\bibitem{options} P. Fran\c{c}ois, G. Gauthier, F. Godin and C. O. P\'erez-Mendoza. Research article, 2025, arXiv:2504.06208v3. \url{https://arxiv.org/abs/2504.06208v3}.
\bibitem{reference} C. O. P\'erez-Mendoza. Author reference implementation and experiment configuration. \url{https://github.com/cpmendoza/deep-hedging_with_options}.
\bibitem{volt} Volt. \emph{Native learner application guide}. Protocol, financial inputs, execution backends and checkpoint support. \url{https://github.com/ingvor-restrat/volt/blob/main/crates/volt-compute/docs-native-learner.md}.
\bibitem{jorm} Jormungandr. \emph{Native learners}. Application endpoints, registered engines and example client. \url{https://github.com/ingvor-restrat/jormungandr/blob/main/docs/native-learners.md}.
\bibitem{authorcode} H. Buehler. Deep Hedging sample implementation. \url{https://github.com/hansbuehler/deephedging}.
\bibitem{burn} Burn. Tensor and automatic-differentiation implementation. \url{https://github.com/tracel-ai/burn}.
\end{thebibliography}
\end{SNcolumns}
\end{document}
